\section{重正化群方程与渐近公式}
\label{sec:4}
\subsection{系数的重正化群方程}
\label{sec:4.1}
我们现在把
\begin{equation}
   \left(\mu\frac{\partial}{\partial\mu}\right)_0
\label{eq:(4.1)}
\end{equation}
作用于式~\eqref{eq:(3.7)} 的两边。由于式~\eqref{eq:(3.7)} 的左边、即式%
~\eqref{eq:(3.5)}，通过流方程~\eqref{eq:(3.1)} 与初始条件~\eqref{eq:(3.4)}
完全由裸量表出，所以 \eqref{eq:(4.1)} 对左边的作用恒等于零。在右边，这一消
没必须对 $t$ 的每个幂次分别成立。于是我们推断：
\begin{align}
   &\left(\mu\frac{\partial}{\partial\mu}\right)_0c_T(t)
   \left\{T_{\mu\nu}\right\}_R(x)=0,
\label{eq:(4.2)}
\\
   &\left(\mu\frac{\partial}{\partial\mu}\right)_0c_S(t)
   \left\{\frac{1}{4}F_{\rho\sigma}^aF_{\rho\sigma}^a\right\}_R(x)=0.
\label{eq:(4.3)}
\end{align}

对于第一个关系~\eqref{eq:(4.2)}，我们回想起能量--动量张量是不重正化的，如
式~\eqref{eq:(2.6)}。于是，用重正化量表示运算~\eqref{eq:(4.1)}，得到
\begin{equation}
   \left(\mu\frac{\partial}{\partial\mu}
   +\beta\frac{\partial}{\partial g}\right)c_T(t)=0.
\label{eq:(4.4)}
\end{equation}

另一方面，对于式~\eqref{eq:(4.3)}，由式~\eqref{eq:(2.11)} 有
\begin{equation}
   \left(\mu\frac{\partial}{\partial\mu}
   +\beta\frac{\partial}{\partial g}
   +\gamma_S\right)c_S(t)=0,
\label{eq:(4.5)}
\end{equation}
其中
\begin{equation}
   \gamma_S
   \equiv-\left(\mu\frac{\partial}{\partial\mu}\right)_0\ln Z_S.
\label{eq:(4.6)}
\end{equation}
式~\eqref{eq:(2.21)}、\eqref{eq:(2.7)} 与~\eqref{eq:(2.17)} 给出
\begin{equation}
   \gamma_S=-g^3\frac{d}{dg}\left(\frac{\beta}{g^3}\right)
   =2b_1g^4+O(g^6).
\label{eq:(4.7)}
\end{equation}

类似地，对式~\eqref{eq:(3.8)}，我们有
\begin{align}
   &\left(\mu\frac{\partial}{\partial\mu}
   +\beta\frac{\partial}{\partial g}\right)
   \left\langle E(t,x)\right\rangle=0,
\label{eq:(4.8)}
\\
   &\left(\mu\frac{\partial}{\partial\mu}
   +\beta\frac{\partial}{\partial g}
   +\gamma_S\right)c_E(t)=0,
\label{eq:(4.9)}
\end{align}
从而
\begin{equation}
   \left(\mu\frac{\partial}{\partial\mu}
   +\beta\frac{\partial}{\partial g}\right)\frac{c_S(t)}{c_E(t)}=0.
\label{eq:(4.10)}
\end{equation}

由标准论证，并基于无量纲量只能通过无量纲组合~$\sqrt{8t}\mu$ 依赖重正化标
度~$\mu$ 这一事实，上述重正化群方程蕴含：
\begin{align}
   &c_T(t)(g;\mu)=c_T(t_0)(\Bar{g}(-\xi);\mu_0),
\label{eq:(4.11)}
\\
   &c_S(t)(g;\mu)
   =\exp\left[\int_0^{-\xi}d\xi'\,\gamma_S\left(\Bar{g}(\xi')\right)\right]
   c_S(t_0)(\Bar{g}(-\xi);\mu_0),
\label{eq:(4.12)}
\\
   &t^2\left\langle E(t,x)\right\rangle(g;\mu)
   =t_0^2\left\langle E(t_0,x)\right\rangle(\Bar{g}(-\xi);\mu_0),
\label{eq:(4.13)}
\\
   &\frac{c_S(t)}{c_E(t)}(g;\mu)=\frac{c_S(t_0)}{c_E(t_0)}(\Bar{g}(-\xi);\mu_0),
\label{eq:(4.14)}
\end{align}
其中对重正化规范耦合与重正化标度的依赖已显式写出。在这些表达式中，跑动耦
合~$\Bar{g}(\xi)$ 定义为
\begin{equation}
   \frac{d\Bar{g}(\xi)}{d\xi}=\beta\left(\Bar{g}(\xi)\right),\qquad
   \Bar{g}(0)=g,
\label{eq:(4.15)}
\end{equation}
并引入变量
\begin{equation}
   \xi\equiv\ln\frac{\sqrt{8t}\mu}{\sqrt{8t_0}\mu_0}.
\label{eq:(4.16)}
\end{equation}

在单圈阶，跑动耦合~\eqref{eq:(4.15)} 由
\begin{equation}
   \Bar{g}(-\xi)^2=\frac{1}{2b_0}\frac{1}{-\xi+1/(2b_0g^2)}
   =\frac{1}{2b_0}\frac{1}{-\ln(\sqrt{8t}\Lambda)+\ln(\sqrt{8t_0}\mu_0)},
\label{eq:(4.17)}
\end{equation}
给出，其中 $\Lambda$ 是单圈水平的 $\Lambda$ 参数，
\begin{equation}
   \Lambda=\mu e^{-1/(2b_0g^2)},
\label{eq:(4.18)}
\end{equation}
而式~\eqref{eq:(4.12)} 中出现的积分为
\begin{equation}
   \int_0^{-\xi}d\xi'\,
   \gamma_S\left(\Bar{g}(\xi')\right)
   =\frac{b_1}{b_0}
   \left[
   g^2-\Bar{g}(-\xi)^2
   \right].
\label{eq:(4.19)}
\end{equation}
在小流时间极限 $t\to0$ 下，$-\xi\to+\infty$，跑动耦
合~$\Bar{g}(-\xi)$~\eqref{eq:(4.17)} 得益于渐近自由而变得非常小。于是，
式~\eqref{eq:(4.11)}--\eqref{eq:(4.14)} 的右端使我们能够用微扰论计算这些系
数的小流时间行为。

\subsection{最低阶近似与渐近公式}
把流方程~\eqref{eq:(3.1)} 的解（见文献~\cite{Luscher:2011bx}）代入式%
~\eqref{eq:(3.7)} 的树图级近似中，我们得到
\begin{equation}
   c_T(t)=g_0^2,
\label{eq:(4.20)}
\end{equation}
理由很简单：我们的能量--动量张量~\eqref{eq:(2.4)} 正比于~$1/g_0^2$。然而，
如果把式~\eqref{eq:(4.17)} 代入并把式~\eqref{eq:(4.11)} 的右端用于这个表达
式，它就会依赖于~$\sqrt{8t_0}\mu_0$，而式~\eqref{eq:(4.11)} 的左端却不依赖
于它。这表明 $c_T(t)$ 应当以一个特定组合的方式依赖~$g^2$
 与~$\sqrt{8t}\mu$，即（对 $D=4$）
\begin{equation}
   c_T(t)=
   g^2\left\{1+2b_0g^2\left[\ln(\sqrt{8t}\mu)+c_1\right]+O(g^4)\right\},
\label{eq:(4.21)}
\end{equation}
其中 $c_1$ 是常数。类似地，由于式~\eqref{eq:(3.7)} 与~\eqref{eq:(3.8)} 中的
最低阶近似给出
\begin{equation}
   c_E(t)=1,\qquad
   c_S(t)=-\frac{b_0}{2}g_0^2\mu^{-2\epsilon},
\label{eq:(4.22)}
\end{equation}
（后者由式~\eqref{eq:(3.10)} 得出），从式~\eqref{eq:(4.14)} 我们有
\begin{equation}
   \frac{c_S(t)}{c_E(t)}
   =-\frac{b_0}{2}
   g^2\left\{1+2b_0g^2\left[\ln(\sqrt{8t}\mu)+c_2\right]+O(g^4)\right\},
\label{eq:(4.23)}
\end{equation}
其中 $c_2$ 是另一个常数。\footnote{把式~\eqref{eq:(4.21)} 用于
式~\eqref{eq:(3.10)}，得
\begin{equation}
   c_S(t)=-\frac{b_0}{2}
   g^2\left\{
   1+2b_0g^2
   \left[
   \ln(\sqrt{8t}\mu)+c_1+\frac{b_1}{2b_0^2}\right]+O(g^4)
   \right\},
\label{eq:(4.24)}
\end{equation}
再用式~\eqref{eq:(4.23)}，得
\begin{equation}
   c_E(t)=1+2b_0g^2
   \left(c_1-c_2+\frac{b_1}{2b_0^2}\right)+O(g^4).
\label{eq:(4.25)}
\end{equation}
}

把式~\eqref{eq:(4.11)} 与~\eqref{eq:(4.14)} 应用于上面的表达式并使用
式~\eqref{eq:(4.17)}，最终得到式~\eqref{eq:(3.11)} 中各系数的渐近行为：
\begin{equation}
   \frac{1}{c_T(t)}
   \stackrel{t\to0+}{\sim}
   -2b_0\left[\ln(\sqrt{8t}\Lambda)+c_1\right]
\label{eq:(4.26)}
\end{equation}
与
\begin{equation}
   \frac{c_S(t)}{c_E(t)}
   \stackrel{t\to0+}{\sim}
   -\frac{b_0}{2}
   \frac{1}{-2b_0\left[\ln(\sqrt{8t}\Lambda)+c_2\right]},
\label{eq:(4.27)}
\end{equation}
从而
\begin{equation}
   \frac{c_S(t)}{c_T(t)c_E(t)}
   \stackrel{t\to0+}{\sim}
   -\frac{b_0}{2}
   \left[1-\frac{c_1-c_2}{-\ln(\sqrt{8t}\Lambda)}\right].
\label{eq:(4.28)}
\end{equation}
也就是说，
\begin{align}
   \left\{T_{\mu\nu}\right\}_R(x)
   &\stackrel{t\to0+}{\sim}
   \biggl\{
   -2b_0\left[\ln(\sqrt{8t}\Lambda)+c_1\right]U_{\mu\nu}(t,x)
\notag\\
   &\qquad\qquad{}
   +\frac{b_0}{2}
   \left[1-\frac{c_1-c_2}{-\ln(\sqrt{8t}\Lambda)}\right]
   \delta_{\mu\nu}
   \left[
   E(t,x)-\left\langle E(t,x)\right\rangle\right]
   \biggr\}.
\label{eq:(4.29)}
\end{align}
这就是我们所寻求的关系：可以从杨--米尔斯梯度流给出的规范不变乘积的小流时
间行为出发，得到正确归一化的守恒能量--动量张量。有趣的是，领头的
$t\to0$ 行为完全不依赖于梯度流的具体定义；其结构与系数仅仅由局域乘积的有
限性和杨--米尔斯理论的可重正化性决定。渐近形式中的次领头修正、即系
数 $c_1$ 与~$c_2$，则依赖于梯度流的具体定义；在附录中我们计算常数 $c_1$
与~$c_2$，得到
\begin{align}
%   c_1&=\ln\sqrt{\pi}+\frac{7}{16}\simeq1.00986,
   c_1&=\ln\sqrt{\pi}+\frac{7}{22}
   \simeq0.890547,
\label{eq:(4.30)}
\\
%   c_2&=\ln\sqrt{\pi}+\frac{3}{44}-\frac{1}{4}
%   +\frac{b_1}{2b_0^2}\simeq0.812034.
   c_2&=\ln\sqrt{\pi}-\frac{7}{44}
   +\frac{b_1}{2b_0^2}
   \simeq0.834762.
\label{eq:(4.31)}
\end{align}

最后，对小流时间确定式~\eqref{eq:(4.29)} 中因子 $\ln(\sqrt{8t}\Lambda)$
（即以单圈 $\Lambda$ 参数~\eqref{eq:(4.18)} 为单位度量流时
间~$t$）的一个可能方法，是使用密度算符的期望值、式~\eqref{eq:(3.6)}。对这个
量，把式~\eqref{eq:(4.13)} 与~\eqref{eq:(4.17)} 应用于单圈计算的结果——文
献~\cite{Luscher:2010iy} 的式 (2.28) 与 (2.29)（特化到纯杨--米尔斯理论）—
—便得到渐近形式
\begin{equation}
   t^2\left\langle E(t,x)\right\rangle
   \stackrel{t\to0+}{\sim}
   \frac{3(N^2-1)}{128\pi^2}
   \frac{1}{-2b_0\left[\ln(\sqrt{8t}\Lambda)+c\right]},
\label{eq:(4.32)}
\end{equation}
其中
\begin{equation}
   c\equiv\ln(2\sqrt{\pi})+\frac{26}{33}-\frac{9}{22}\ln3
   \simeq1.60396.
\label{eq:(4.33)}
\end{equation}
可以把这个渐近表示用作式~\eqref{eq:(4.29)} 中 $\ln(\sqrt{8t}\Lambda)$
 的取值。\footnote{实际操作中，人们会用式~\eqref{eq:(4.29)} 从%
~$t^2U_{\mu\nu}(t,x)$ 与~$t^2E(t,x)$ 计算 $t^2\{T_{\mu\nu}\}_R(x)$。然后由
$\sqrt{8t}\Lambda$ 的值可以推出 $\{T_{\mu\nu}\}_R(x)/\Lambda^4$。}
